\themaN
\graphicspath{{../../S08_Expressions_algebriques/Images/}}

\chapter{Expressions\\algébriques}
\label{S08}


%%%%%%%%%%%%%%%%%%%%%%%%%%%%%%
%%%%%%%%%%%%%%%%%%%%%%%%%%%%%%
\begin{autoeval}
   \small
   \begin{enumerate}
      \item Il utilise les notations $2a$ pour $a\times2$ ou $2\times a$ et $ab$ pour $a\times b$, $a^2$ pour $a\times a$ et $a^3$ pour $a\times a\times a$.
      \item Il produit une expression littérale pour élaborer une formule ou traduire un programme de calcul.
      \item Il utilise une lettre pour traduire des propriétés générales.
      \item Il utilise une lettre pour démontrer une propriété générale.
      \item Il substitue une valeur numérique à une lettre pour : calculer la valeur d’une expression littérale ; tester, à la main ou de façon instrumentée, si une égalité où figurent une ou deux indéterminées est vraie quand on leur attribue des valeurs numériques ; contrôler son résultat.
   \end{enumerate}
\end{autoeval}

\begin{prerequis}
   \begin{itemize}
      \item Notion d'inconnue.
      \item[\com] Utiliser le calcul littéral pour traduire une propriété générale, pour démontrer un résultat général, pour valider ou réfuter une conjecture, pour modéliser une situation.
   \end{itemize}
\end{prerequis}

\vfill

\begin{debat}[Débat : petits contes mathématiques]
   Les vidéos des {\bf Petits contes mathématiques} ont été créés par {\it Universciences} : c'est une série pédagogique qui retrace l'histoire des maths à travers la découverte d'une notion, d'une formule, d'une conjecture ou d'une équation. Le récit est rythmé par des illustrations animées et la légèreté du ton dédramatise le sujet pour tous ceux qui ne seraient pas des \og matheux \fg. \\
   \begin{center}
      \begin{pspicture}(0,0)(5,3)
         \textcolor{B1}{
         \rput{10}(1,1){\huge $x$}
         \rput{20}(3.5,2.25){\huge $y$}
         \rput{-10}(3,0.5){\huge $z$}
         \rput(0.5,2.5){\huge $+$}
         \rput(2,1.5){\huge $\times$}
         \rput(4.5,2){\huge $()$}
         \rput(5,0.75){\huge $-$}
         \rput(2,3){\huge $\div$}
         \rput{-20}(5,3){\huge $t$}}
      \end{pspicture}
   \end{center}
   \bigskip
   \begin{cadre}[B2][J4]
      \begin{center}
         Vidéo : \href{https://leblob.fr/fondamental/le-x}{\bf Le $x$}, site Internet {\it Le blob}, épisode de la série {\it Petits contes mathématiques}.
      \end{center}
   \end{cadre}
\end{debat}


%%%%%%%%%%%%%%%%%%%%%%%%%%%%%%%%
%%%%%%%%%%%%%%%%%%%%%%%%%%%%%%%%
\activites

\begin{activite}[Des variables aux inconnues]
   {\bf Objectifs :} voir l'effet des variables dans un programme de calcul créé avec Scratch ; produire une expression littérale.
   \begin{QCM}
   On considère le programme de calcul ci-dessous dans lequel $x$, Etape 1, Etape 2 et Résultat sont quatre variables.
   \begin{center}
      \begin{Scratch}[Echelle=0.8]
         Place Drapeau ;
         Place Demander("Choisis un nombre");
         Place MettreVar("x",OvalCap("réponse"));
         Place DireT("Je multiplie le nombre par 6.","2");
         Place MettreVar("Etape 1",OpMul("6",OvalVar("x")));
         Place DireT("J'ajoute 10 au résultat.","2");
         Place MettreVar("Etape 2",OpAdd(OvalVar("Etape 1"),"10"));
         Place DireT("Je divise le résultat par 2.","2");
         Place MettreVar("Résultat",OpDiv(OvalVar("Etape 2"),"2"));
         Place Dire(OpRegrouper("J'obtiens finalement",OvalVar("Résultat")));
      \end{Scratch}
   \end{center}
   \begin{enumerate}
   {\psset{unit=0.9}
      \item Julien a fait fonctionner ce programme en choisissant le nombre 5.\\
      Vérifier que ce qui est dit à la fin est : \og J’obtiens finalement 20 \fg. Pour cela, remplir le diagramme suivant : \\
         \begin{pspicture}(-4,-0.2)(8,1.8)
            \psframe(0,0.5)(1,1.5)
            \rput(0.5,0){$x$}
            \psline[arrowsize=2mm]{->}(1.1,1)(2.4,1)
            \psframe(2.5,0.5)(3.5,1.5)
            \rput(3,0){étape 1}
            \psline[arrowsize=2mm]{->}(3.6,1)(4.9,1)
            \psframe(5,0.5)(6,1.5)
            \rput(5.5,0){étape 2}
            \psline[arrowsize=2mm]{->}(6.1,1)(7.4,1)
            \psframe(7.5,0.5)(8.5,1.5)
            \rput(8,0){résultat}
         \end{pspicture}         
      \item Que dit le programme si Zakariae le fait fonctionner en choisissant au départ le nombre 7 ? \\
          \begin{pspicture}(-4,-0.2)(8,1.8)
            \psframe(0,0.5)(1,1.5)
            \rput(0.5,0){$x$}
            \psline[arrowsize=2mm]{->}(1.1,1)(2.4,1)
            \psframe(2.5,0.5)(3.5,1.5)
            \rput(3,0){étape 1}
            \psline[arrowsize=2mm]{->}(3.6,1)(4.9,1)
            \psframe(5,0.5)(6,1.5)
            \rput(5.5,0){étape 2}
            \psline[arrowsize=2mm]{->}(6.1,1)(7.4,1)
            \psframe(7.5,0.5)(8.5,1.5)
            \rput(8,0){résultat}
         \end{pspicture}         
      \item Titouan fait fonctionner le programme, et ce qui est dit à la fin est : \og J’obtiens finalement 8 \fg. Quel nombre a-t-il choisi au départ ? \\
      \begin{pspicture}(-4,-0.2)(8,1.8)
            \psframe(0,0.5)(1,1.5)
            \rput(0.5,0){$x$}
            \psline[arrowsize=2mm]{<-}(1.1,1)(2.4,1)
            \psframe(2.5,0.5)(3.5,1.5)
            \rput(3,0){étape 1}
            \psline[arrowsize=2mm]{<-}(3.6,1)(4.9,1)
            \psframe(5,0.5)(6,1.5)
            \rput(5.5,0){étape 2}
            \psline[arrowsize=2mm]{<-}(6.1,1)(7.4,1)
            \psframe(7.5,0.5)(8.5,1.5)
            \rput(8,0){résultat}
         \end{pspicture} 
      \item Si l’on appelle $x$ le nombre choisi au départ, écrire en fonction de $x$ l’expression obtenue à la fin du programme. \par \medskip
         \pointilles}
   \end{enumerate}
   \end{QCM}
   \vfill\hfill{\footnotesize\it Cette activité est une adaptation du DNB 2017 de Pondichery.}
\end{activite} 


%%%%%%%%%%%%%%%%%%%%%%%%%%%%%
%%%%%%%%%%%%%%%%%%%%%%%%%%%%%
\cours 

\section{Expressions littérales} %%%%%

\medskip

\begin{definition}
   Une {\bf expression littérale} est une succession d'opérations où apparaissent des lettres représentant des nombres. On parle aussi d'{\bf expression algébrique}.
\end{definition}

\begin{exemple}
   $4\times x+3$ \; ; \; $2\times a-89$ \; ; \; $y-3\times z+t-4$.
\end{exemple}

\bigskip

Lorsque l'on choisit une certaine valeur pour chaque lettre d'une expression littérale, on peut calculer la valeur de l'expression littérale.  

\begin{exemple}[0.5]
   Calculer $4\times x+3\times y$ pour $x=5$ et $y =0$. \\
   \correction
      $4\times{\blue x}+3\times{\red y} =4\times{\red 5}+3\times{\blue0} =20+0 =20$.
\end{exemple}


%%%%%%%%%%%%%%%%%%%%%
\section{Produire une expression littérale}

Parfois, on est amené à trouver une expression littérale pour avoir une formule générale.

\begin{exemple}[0.55]
\ \\ [-10mm]
   \begin{itemize}
      \item Matéo a quatre ans de plus que Noé. \\
         Exprimer l'âge de Matéo par rapport à celui de Noé.
      \item Un rectangle a pour largeur $\ell$ et pour longueur $L$. \\
         Donner l'expression de son aire et de son périmètre.
   \end{itemize}
   \correction
      \ \\ [-10mm]
      \begin{itemize}
         \item En appelant \og $\red x$ \fg{} l'âge de Noé, l'âge de Matéo peut s'écrire \og ${\red x}+4$ \fg.
         \item $\mathcal{A} =\ell\times L$. \\
         $\mathcal{P} =2\times(\ell+L)$.
      \end{itemize}
\end{exemple}

\medskip


%%%%%%%%%%%%%%%%%%%%%%%%%%%%%%%
\section{Simplification d'une expression littérale}

Pour simplifier une expression littérale, on peut supprimer le signe \og $\times$ \fg{} devant une lettre, une parenthèse ou entre deux lettres. 

\begin{propriete}
   \begin{itemize}
      \item On utilise la notation $2a$ pour $a+a$ ou $a\times2$ ou encore $2\times a$. On dit \og deux a \fg.
      \item On utilise la notation $ab$ pour $a\times b$. On dit \og ab \fg.
      \item On utilise la notation $a^2$ pour $a\times a$. On dit \og a au carré \fg.
      \item On utilise la notation $a^3$ pour $a\times a\times a$. On dit \og a au cube \fg. \\ [-8mm]
   \end{itemize}
\end{propriete}

\bigskip

On écrit une expression comportant un nombre et une \og lettre \fg{} avec le nombre précédé de la \og lettre \fg.

\begin{exemple}[0.5]
   $A =5\times y-2\times x$. \\
   $B =x\times x \times x + 7\times y\times y$. \\
   $C =x\times3$
   \correction
      $A =5y-2x$. \\
      $B =x^3+7y^2$. \\
      $C =3x$, est préférable à $C=x3$. Par commutativité de la multiplication, $x\times3 =3\times x =3x$.
\end{exemple}


%%%%%%%%%%%%%%%%%%%%%%%%%%%%%%%%
%%%%%%%%%%%%%%%%%%%%%%%%%%%%%%%%
\exercicesbase

\begin{colonne*exercice}

\begin{exercice} %1
   Calculer les expressions $A$ à $G$ pour $x =2$ et $y =0$.
   \begin{enumerate}
      \item $A =2x$
      \item $B =7-3y$
      \item $C =4x+5$
      \item $D =4(x-2)(x+3)$
      \item $E =2xy$
      \item $F =(x-y)\times3$
      \item $G =2y^2+3x-5$
   \end{enumerate}
\end{exercice}

\begin{corrige}
   \ \\ [-5mm]
   \begin{enumerate}
      \item $A =2x =2\times2 =\blue4$
      \item $B =7-3y =7-3\times0 =\blue7$
      \item $C =4x+5 =4\times2+5 =\blue13$
      \item $D =4(x-2)(x+3) =4(2-2)(2+3)$ \\
         \hspace*{7.6mm} $=4\times0\times5 =\blue0$
      \item $E =2xy =2\times2\times0 =\blue0$
      \item $F =(x-y)\times3 =(2-0)\times3 =2\times3 =\blue6$
      \item $G =2y^2+3x-5 =2\times0^2+3\times2-5$ \\
         \hspace*{7.4mm} $=6-5 =\blue1$
   \end{enumerate}
\end{corrige}

\bigskip


\begin{exercice} %2
   Compléter le tableau suivant. \\ [2mm]
   {\hautab{1.5}
   \begin{Ctableau}{\linewidth}{8}{c}
      \hline
      $x$ & $0$ & $1$ & $2$ & $3$ & $4$ & $5$ & $6$ \\
      \hline
      $2x+4$ & & & & & & & \\
      \hline
      $x^2+1$ & & & & & & & \\
      \hline
   \end{Ctableau}}
\end{exercice}

\begin{corrige}
   \ \\ [-3mm]
   \begin{Ctableau}{\linewidth}{8}{c}
      \hline
      $x$ & $0$ & $1$ & $2$ & $3$ & $4$ & $5$ & $6$ \\
      \hline
      $2x+4$ & \textcolor{blue}{4} & \textcolor{blue}{6} & \textcolor{blue}{8} & \textcolor{blue}{10} & \textcolor{blue}{12} & \textcolor{blue}{14} & \textcolor{blue}{16} \\
      \hline
      $x^2+1$ & \textcolor{blue}{1} & \textcolor{blue}{2} & \textcolor{blue}{5} & \textcolor{blue}{10} & \textcolor{blue}{17} & \textcolor{blue}{26} & \textcolor{blue}{37} \\
      \hline
   \end{Ctableau}
\end{corrige}

\bigskip


\begin{exercice} %3
   Si $x$ représente un nombre, comment peut-on écrire les expressions suivantes :
   \begin{enumerate}
      \item Le double de $x$.
      \item Le tiers de $x$.
      \item La somme de $x$ et de 13.
      \item La différence de $x$ et de 7.
      \item Le triple de la somme de 2 et de $x$.
      \item Le tiers de la différence de 16 et $x$.
   \end{enumerate}
\end{exercice}

\begin{corrige}
   \ \\ [-5mm]
   \begin{enumerate}
      \item $2\times x =\blue2x$ \smallskip
      \item $\dfrac13\times x =\blue\dfrac13x =\dfrac{x}{3}$ \smallskip
      \item $\blue x+13$
      \item $\blue x-7$
      \item $3\times(2+x) =\blue3(2+x)$
      \item $\dfrac13\times(16-x) =\blue\dfrac13(16-x) =\dfrac{16-x}{3}$
   \end{enumerate}
\end{corrige}

\bigskip


\begin{exercice} %4
   Si on note $z$ l'âge en années de Rose aujourd'hui, comment peut-on noter :
   \begin{enumerate}
      \item L'âge qu'elle aura dans deux ans ?
      \item Le triple de l'âge qu'elle avait il y a quatre ans ?
      \item La moitié de l'âge qu'elle aura dans cinq ans ?
      \item Son année de naissance si on est en 2022 ?
   \end{enumerate}
\end{exercice}

\begin{corrige}
   \ \\ [-5mm]
   \begin{colenumerate}{2}
      \item $\blue z+2$
      \item $\blue 3(z-4)$
      \item $\blue \dfrac12(z+5) =\dfrac{z+5}{2}$ \smallskip
      \item $\blue 2021-z$
   \end{colenumerate}
\end{corrige}

\bigskip


\begin{exercice} %5
   Voici un programme :
   \begin{center}
      \ProgCalcul[Enonce,Largeur=4.5cm]{
         Choisis un nombre,
         Retire-lui 5,
         Multiplie le résultat par 3}
   \end{center}
   \begin{enumerate}
      \item Faire fonctionner le programme avec trois nombres de son choix supérieurs ou égaux à 5.
      \item Quel nombre faut-il choisir pour obtenir 6 ?
      \item Soit $x$ le nombre de départ, donner l'expression finale en fonction de $x$.
   \end{enumerate}
\end{exercice}

\begin{corrige}
   \ \\ [-5mm]
   \begin{enumerate}
      \item \ProgCalcul{10,-5 *3}. Avec 0, on obtient {\blue 15}. \\ [1mm]
         \ProgCalcul{12,-5 *3}. Avec 12, on obtient {\blue 21}. \\ [1mm]
         \ProgCalcul{5,-5 *3}. Avec 5 on obtient {\blue 0}. \smallskip
      \item \ProgCalcul[Direct=false]{7,-5 *3}. On obtient 6 avec {\blue 7}. \smallskip
      \item $x\xrightarrow{-5}x-5\xrightarrow{\times3}\blue3(x-5)$.
   \end{enumerate}
\end{corrige}

\bigskip


\begin{exercice} %6
   Sahya construit un carré $ABCD$ de \ucm{5} de côté.
   \begin{enumerate}
      \item Calculer le périmètre $\mathcal{P}_1$ et l'aire $\mathcal{A}_1$ de $ABCD$.
      \item On augmente ses côtés de $k$ cm. Exprimer, en fonction de $k$ :
      \begin{itemize}
         \item la longueur $L$ du nouveau côté ;
         \item le nouveau périmètre $\mathcal{P}_2$ de ce carré ;
         \item la nouvelle aire $\mathcal{A}_2$ de ce carré.
      \end{itemize}
      \item Grâce aux expressions trouvées en 2, donner le périmètre et l'aire si on augmente le côté de 2 cm.
   \end{enumerate}
   {\psset{unit=1.2}
   \begin{pspicture}(-2,-1)(3,3.5)
      \pstGeonode[PosAngle={-135,-45,45,135},PointSymbol=none](0,0){A}(2,0){B}(2,2){C}(0,2){D}
      \psframe(0,0)(2,2)
      \psframe(0,0)(0.2,0.2)
      \psframe(0,0)(3,3)
      \psline{<->}(0,-0.5)(2,-0.5)
      \rput(1,-0.75){\small \ucm{5}}
      \psline{<->}(2,-0.5)(3,-0.5)
      \rput(2.5,-0.75){\small k}
   \end{pspicture}}
\end{exercice}

\begin{corrige}
   \ \\ [-5mm]
   \begin{enumerate}
      \item $\mathcal{P}_1 =4\times\ucm{5} =\blue\ucm{20}$ ; $\mathcal{A}_1 =(\ucm{5})^2 =\blue\ucmq{25}$
      \item Les longueurs sont en \ucm{} et les aires en \ucmq{}.
      \begin{itemize}
         \item La longueur $L$ du nouveau côté est $\blue5+k$.
         \item Le nouveau périmètre vaut $\mathcal{P}_2 =\blue4(5+k)$
         \item La nouvelle aire vaut $\mathcal{A}_2 =\blue(5+k)^2$
      \end{itemize}
      \item $\mathcal{P}_2 =4(5+2) \, \ucm{} = 4\times7 \, \ucm{} =\blue\ucm{28}$ et \\
      \quad\, $\mathcal{A}_2 =((5+2)\,\ucm{})^2 =(\ucm{7})^2 =\blue\ucmq{49}$
   \end{enumerate}
\end{corrige}

\bigskip


\begin{exercice} %7
Recopier les expressions suivantes en supprimant les signes $\times$ s'ils ne sont pas nécessaires.
\begin{enumerate}
   \item $A =9\times n$
   \item $B =x\times3$
   \item $C =12\times(a-3)$
   \item $D =2\times a(2\times8)$
   \item $E =n\times x$
   \item $F =2\times\pi\times R$
\end{enumerate}
\end{exercice}

\begin{corrige}
   \ \\ [-5mm]
   \begin{enumerate}
      \item $A =9\times n =\blue9n$
      \item $B =x\times3 =\blue3x$
      \item $C =12\times(a-3) =\blue12(a-3)$
      \item $D =2\times a(2\times8) =2\times16a =\blue32a$
      \item $E =n\times x =\blue nx$
      \item $F =2\times\pi\times R =\blue2\pi R$
\end{enumerate}
\end{corrige}

\bigskip

\begin{exercice} %8
   Simplifier les expressions suivantes :
   \begin{enumerate}
      \item $A =a\times a$
      \item $B =b\times b\times b$
      \item $C =c\times c\times3$
      \item $D =9+d\times d\times d$
      \item $E =a\times a\times b\times3$
      \item $F =x\times x\times x-2\times y\times y$
      \item $G =(a+b)\times(a+b)$
      \item $H =(x+y)(x+y)(x+y)$
   \end{enumerate}
\end{exercice}

\begin{corrige}
   \ \\ [-5mm]
   \begin{enumerate}
      \item $A =a\times a =\blue a^2$
      \item $B =b\times b\times b =\blue b^3$
      \item $C =c\times c\times3 =\blue3c^2$
      \item $D =9+d\times d\times d =\blue9+d^3$
      \item $E =a\times a\times b\times3 =\blue3a^2b$
      \item $F =x\times x\times x-2\times y\times y =\blue x^3-2y^2$
      \item $G =(a+b)\times(a+b) =\blue(a+b)^2$
      \item $H =(x+y)(x+y)(x+y) =\blue(x+y)^3$ \medskip
   \end{enumerate}
\end{corrige}

\flushright{\footnotesize\it D'après Sésamath, le manuel de cycle 4. Magnard 2016}

\end{colonne*exercice}


%%%%%%%%%%%%%%%%%%%%%%%%%%%%%%%
%%%%%%%%%%%%%%%%%%%%%%%%%%%%%%%
\Recreation

\def\croix{\psframe(0,-1)(1,2) \psframe(-1,-0)(2,1)}
\def\carre{\psframe(0,0)(1,1)}

\begin{enigme}[Défis !!!]

\vfill

\parbox{1.75cm}{\includegraphics[width=2cm]{defi} \\ [2mm] \textcolor{orange}{\bf\large \, Défi 1}}
\qquad
\parbox{14.45cm}{
   Rayan et Aya ont choisi un nombre (entier positif). \\
   Rayan le multiplie par 5 et ajoute 35. \\
   Aya le multiplie par 2 et ajoute 146. \\
   Ils trouvent le même nombre à la fin. \\
   Quel nombre ont-ils choisi ?}

\vfill

\parbox{1.75cm}{\includegraphics[width=2cm]{defi} \\ [2mm] \textcolor{orange}{\bf\large \, Défi 2}}
\qquad
\parbox{14cm}{
   Avec des petits carrés tous identiques, on construit un pattern selon le modèle évolutif ci-dessous : \\
   {\psset{unit=0.8}
   \begin{pspicture}(-3.5,-4.5)(4,3.75)
      \croix
      \rput(0.5,-2){Rang 1}
   \end{pspicture}
   \begin{pspicture}(-5,-4.5)(5,3.75)
      \croix
      \rput(2,0){\carre}
      \rput(0,2){\carre}
      \rput(-2,0){\carre}
      \rput(0,-2){\carre}
      \rput(0.5,-3){Rang 2}
   \end{pspicture}
   
   \begin{pspicture}(-8,-4.5)(4,1)
      \croix
      \rput(2,0){\carre}
      \rput(3,0){\carre}
      \rput(0,2){\carre}
      \rput(0,3){\carre}
      \rput(-2,0){\carre}
      \rput(-3,0){\carre}
      \rput(0,-2){\carre}
      \rput(0,-3){\carre}
      \rput(0.5,-4){Rang 3}
   \end{pspicture}}
   \begin{enumerate}
      \item Dessiner l’élément du rang suivant et expliquer la règle.
      \item Déterminer le nombre de petits carrés des éléments du rang 5, du rang 10, du rang 17.
      \item Déterminer le nombre de petits carrés de l’élément du rang 100 et donner un moyen de calculer rapidement le nombre de petits carrés d’un élément à n’importe quel rang.
     \item Existe-t-il un élément qui contient 532 petits carrés ? Un élément qui contient 813 petits carrés ?
  \end{enumerate}}
   
   \vfill\hfill {\it\small Source : La résolution de problèmes mathématiques au collège, MENJS, 2021}
 \end{enigme}  
 
\begin{corrige}
\hspace*{-7.5mm} \textcolor{orange}{\bf Défi 1} \\
   On peut modéliser ce défi par un schéma en barres. Soit $x$, le nombre choisi. \\
   \ModeleBarre[Separation=2]{DeepSkyBlue 1 "$x$" 1 "$x$" 1 "$x$" 1 "$x$" 1 "$x$" PowderBlue 2 "35"}{DeepSkyBlue 1 "$x$" 1 "$x$" SkyBlue 5 "146"}
   \ModeleBarre[Separation=3]{DeepSkyBlue 1 "$x$" 1 "$x$" 1 "$x$" PowderBlue 2 "35"}{ SkyBlue 3 "111" PowderBlue 2 "35"}
   \ModeleBarre{DeepSkyBlue 1 "$x$" 1 "$x$" 1 "$x$"}{SkyBlue 3 "111"} donc, $x =111\div3 ={\blue 37}$.

\hspace*{-7.5mm} \textcolor{orange}{\bf Défi 2} \\
   \begin{enumerate}
      \item Pour le rang 4, il suffit {\blue d'ajouter 1 carré à chacune des extrémités}.
      \item Au rang 1, on a 5 carrés. Au rang 2, $5+4 =9$ carrés. Au rang 3, $9+4 =13$ carrés.  Au rang 4, $13+4 =17$ carrés. {\blue Au rang 5, $17+4 =21$ carrés}. \\
         En continuant ainsi, on aura {\blue 41 carrés au rang 10 et 69 carrés au rang 17}.
         \item Soit $n$ le rang demandé, {\blue le nombre de carrés au rang $n$ est de $1+4\times n$}. \\
            Au rang 100, cela fait donc $1+4\times100 =401$.
         \item En enlevant 1 carré, on doit obtenir un multiple de 4. Or, $532-1 =531$ n'est pas un multiple de 4. \\
            $813-1 =812 =4\times203$ est un multiple de 4. \\
            {\blue Il y a donc 813 carrés au rang 203}.     
   \end{enumerate}
\end{corrige}

